\documentclass[
    pdftex,
    %draft, % Prints a box at the end of overfull lines, changes behavior of some packages
    final,
    11pt,
    a4paper,
    parskip=false, % Paragraphs are identified by indentation of the first line by 1em
    twoside, % Prepares the document for two-sided printing
    footheight=0mm, % remove footer
    footinclude=false,
    toc=bibliography, % add bibliography to table of contents
    toc=listof % add list of figures and tables to table of contents
]{scrbook} % documentation: https://mirrors.ctan.org/macros/latex/contrib/koma-script/doc/scrguide-en.pdf

% language and encoding
\usepackage[utf8]{inputenc} % set document encoding to UTF-8
\usepackage[main=english, ngerman]{babel} % adjust language depending on thesis language (affects hyphenation)
\usepackage[T1]{fontenc} % set font encoding so that special characters are displayed correctly

% configure graphics
\usepackage[pdftex]{graphicx}

% fonts
%\usepackage{mathpazo} % Palatino font (serif)
%\usepackage{fourier} % Utopia font (serif)
\usepackage{lmodern} % Latin Modern Roman (helvet and courier below replace its sans-serif and typewriter fonts)
\usepackage[scaled=0.9]{helvet} % URW Nimbus Sans (Helvetica clone)
\usepackage{courier} % Courier Font (typewriter)

% math formulas
\usepackage{amsmath}
\usepackage{amssymb} 
\usepackage{eucal} % more curly versions for \mathcal{...}
\usepackage{nicefrac} % nicer fractions
\usepackage{bm} % bold math symbols

% outlined characters
\usepackage{contour}
\contourlength{0.1pt}

% scientific notation
\usepackage{siunitx}
\sisetup{output-exponent-marker=\ensuremath{\mathrm{e}}}

% page margins
\usepackage[paper=a4paper, inner=30mm, outer=25mm, top=30mm, bottom=25mm, head=21pt, heightrounded]{geometry} % heightrounded: the text height is a whole number of lines

% references
\usepackage[square, comma, sort, numbers]{natbib} % sort&compress

% citation commands (not compatible with natbib)
%\usepackage{cite}

% prevent single lines at the start of a paragraph (Schusterjungen)
\clubpenalty = 10000
% prevent single lines at the end of a paragraph (Hurenkinder)
\widowpenalty = 10000 \displaywidowpenalty = 10000
% end a page early instead of stretching the space between paragraphs when a single line would remain
\raggedbottom

% enables text wrapping around figures.
\usepackage{wrapfig} 

% allows you to insert \FloatBarrier commands at places that floats should not move past, perhaps at every \section
\usepackage{placeins}

% 1st, 2nd, 3rd, etc.
\usepackage[super]{nth}

% unnumbered theorem-like environment for definitions
\usepackage{amsthm}
\newtheoremstyle{definitiontight}% name
  {.5\baselineskip plus .2\baselineskip minus .2\baselineskip}% Space above
  {.6\baselineskip plus .2\baselineskip minus .2\baselineskip}% Space below
  {}% Body font
  {}%Indent amount (empty = no indent, \parindent = para indent)
  {\bfseries}%  Thm head font
  {.}%       Punctuation after thm head
  {.5em}%      Space after thm head: " " = normal interword space;
        %       \newline = linebreak
  {}%       Thm head spec
\theoremstyle{definitiontight}
\newtheorem{definition}{Definition}[section]
\newtheorem{hypothesis}{Hypothesis}[section]
% https://tex.stackexchange.com/a/347781/100364
\renewcommand{\thehypothesis}{\arabic{hypothesis}}% Remove section from theorem counter representation

% colors for boxes, listings, links, and headers
\usepackage{xcolor} % load before \definecolor so that the colors are compatible with packages that use xcolor
\definecolor{light-gray}{gray}{0.97}
\definecolor{gray}{rgb}{0.4,0.4,0.4}
\definecolor{darkblue}{rgb}{0.0,0.0,0.6}
\definecolor{cyan}{rgb}{0.0,0.6,0.6}
\definecolor{background-gray}{gray}{0.96}

% configure first page of each chapter (thanks to Fabian Beck, Rainer Lutz, Benjamin Biegel)
\makeatletter% see the German TeX FAQ (De-TeX-FAQ)
\renewcommand*{\chapterformat}{
  \begingroup% keeps the change of \unitlength local
    \setlength{\unitlength}{1mm}%
    \begin{picture}(20,20)(0,3)
      \setlength{\fboxsep}{0pt}
       \put(0,0){\makebox(20,20){\rule{20\unitlength}{20\unitlength}}}%
       %\put(0,21){\makebox(20,20){\rule{20\unitlength}{20\unitlength}}}
      % horizontal rule
      \put(22,13){\linethickness{0.5pt}\line(1,0){\dimexpr
          \textwidth-24\unitlength\relax\@gobble}}
      % chapter number
      \put(-1.4,0.8){\makebox(18,18)[r]{%
          \color{white}\fontsize{22\unitlength}{28\unitlength}\selectfont\thechapter
          %\kern-.05em% shifts the number to the right within its box
        }}
      % label "Chapter N"
      \put(20,13){\makebox(\dimexpr
          \textwidth-18\unitlength\relax\@gobble,\ht\strutbox\@gobble)[l]{%
            \ \normalsize\color{black}\chapapp~\thechapter
          }}
    \end{picture} % <-- the space before this comment is intentional
  \endgroup
}
\makeatother

% fancy headings (thanks to Fabian Beck, Rainer Lutz, Benjamin Biegel, and ChatGPT)
\usepackage{scrlayer-scrpage}
\usepackage[listings=false]{scrhack} % suppress warnings
\pagestyle{scrheadings}
\clearpairofpagestyles 
\lehead{\colorbox{black}{\framebox[1cm][c]{\textcolor{white}{\textbf{\small\thepage}}}}\hspace{0.4cm}{\leftmark}}
\rohead{\rightmark\hspace{0.4cm}\colorbox{black}{\framebox[1cm][c]{\textcolor{white}{\textbf{\small\thepage}}}}}
\setkomafont{pageheadfoot}{\fontfamily{pag}\selectfont\footnotesize}
\setlength{\footheight}{\baselineskip}
\KOMAoptions{headsepline=0.5pt}

% configure environment for quotations to be used at the beginning of a chapter (dictum)
\newcommand{\quotemarks}[1]{{\bfseries{#1}}}
\newenvironment{chapter-quotation}[1]
{\begin{quotation}\pushQED{#1}\noindent\slshape\quotemarks{``}\ignorespaces}
{\unskip\quotemarks{''}\par\medskip\noindent\hspace*{\fill}{\normalfont\fontsize{10}{12}\selectfont---\popQED}\end{quotation}}
%\newenvironment{chapter-quotation}[1]
%{\pushQED{#1}\begin{quotation}``''}
%{\begin{flushright}---\popQED\end{flushright}\end{quotation}}

% tables
\usepackage{multirow}

% editorial commands
\newcommand{\boxedtext}[1]{\fbox{\scriptsize\bfseries\textsf{#1}}}
\newcommand{\nota}[2]{%
	\boxedtext{#1}
	{\small$\blacktriangleright$\emph{\textsl{#2}}$\blacktriangleleft$}%
}
\newcommand{\todo}[1]{{\color{red}\nota{TODO}{#1}}} % important: Remove spacing before and after \todo{}
\newcommand{\remark}[2]{{\color{blue}\nota{#1}{#2}}} % important: Remove spacing before and after \remark{}
% uncomment the following line to hide all \todo{} and \remark{} commands
%\renewcommand\nota[2]{}

% customize links and references
\usepackage[hyphens]{url} % line breaks in URLs
\usepackage[colorlinks=true, linkcolor=darkblue, citecolor=darkblue, urlcolor=darkblue]{hyperref} % dark blue links, readable on screen and in print
\renewcommand{\chapterautorefname}{Chapter}
\renewcommand{\sectionautorefname}{Section}
\renewcommand{\subsectionautorefname}{Section}
\renewcommand{\subsubsectionautorefname}{Section}
% babel overrides the labels...
\addto\extrasenglish{%
  \renewcommand{\chapterautorefname}{Chapter}
  \renewcommand{\sectionautorefname}{Section}
  \renewcommand{\subsectionautorefname}{Section}
  \renewcommand{\subsubsectionautorefname}{Section}
}

% configure enumerate, itemize, and description environments
\usepackage{enumitem}

% arithmetic with lengths, e.g., \widthof{<text>} for the width of a text
\usepackage{calc}

% boxes for results
\usepackage[framemethod=tikz]{mdframed}
\newmdenv[
  innerleftmargin=7pt,
  innerrightmargin=7pt,
  skipabove=\baselineskip,
  skipbelow=\baselineskip,
  tikzsetting={draw=black,dashed,line width=0.5pt,dash pattern = on 4pt off 2pt},
  linecolor=white,
  backgroundcolor=white
]{dashedbox}
\newmdenv[
  innerleftmargin=7pt,
  innerrightmargin=7pt,
  skipabove=\baselineskip,
  skipbelow=\baselineskip,
  tikzsetting={draw=black, line width=0.5pt},
  linecolor=black,
  backgroundcolor=white
]{normalbox}
\newmdenv[
  topline=false,
  bottomline=false,
  rightline=false,
  skipabove=\topsep,
  skipbelow=\topsep,
  innertopmargin=0pt,
  innerbottommargin=0pt,
  innerleftmargin=7pt,
  innerrightmargin=0pt,
  tikzsetting={draw=black, line width=3pt},
  linecolor=black,
  backgroundcolor=white
]{verticalline}
\newmdenv[
  topline=false,
  bottomline=false,
  rightline=false,
  skipabove=\baselineskip,
  skipbelow=\baselineskip,
  innertopmargin=6pt,
  innerbottommargin=6pt,
  innerleftmargin=12pt,
  innerrightmargin=12pt,
  tikzsetting={draw=black, line width=4pt},
  linecolor=black,
  backgroundcolor=background-gray
]{verticalline-background}

% source code listings
\usepackage{listings}

\lstset{
  basicstyle=\ttfamily\small,
  showspaces=false,
  showstringspaces=false,
  showtabs=false,
  tabsize=2,
  breaklines=true,
  breakatwhitespace=true,
  aboveskip=0.5\baselineskip,
  belowskip=0.5\baselineskip,
  xleftmargin=0.5\parindent,
  xrightmargin=0pt,
  keywordstyle=\color{blue},
  identifierstyle=\color{black},
  commentstyle=\color{gray},
  stringstyle=\color{gray},
  emphstyle=\color{red},
}

\lstnewenvironment{plain} {\lstset{
}}{}

\lstnewenvironment{java} {\lstset{
  language=Java
}}{}

\lstnewenvironment{java-comment} {\lstset{
  language=Java,
  commentstyle=\color{black}
}}{}

\lstnewenvironment{xml} {\lstset{
  language=XML
}}{}

\lstnewenvironment{code} {\lstset{
}}{}

\lstnewenvironment{regex} {\lstset{
  alsoletter=\\()\{\}.,
  morekeywords={String, long, if, return, format, B, int, Math, log}
}}{}

% gray background for table cells
\usepackage{colortbl} % provides \cellcolor
\definecolor{VeryLightGray}{rgb}{0.92,0.92,0.92}
\newcommand{\graybg}{\cellcolor{VeryLightGray}}

% dummy text
\usepackage{lipsum}

% configure captions
\usepackage[font=small, labelfont=bf, labelsep=endash, format=plain]{caption}

% tables
\usepackage{tabularx}
\usepackage{booktabs}

% custom table column types (http://tex.stackexchange.com/questions/12703/how-to-create-fixed-width-table-columns-with-text-raggedright-centered-raggedlef)
\usepackage{array}
\newcolumntype{L}[1]{>{\raggedright\let\newline\\\arraybackslash\hspace{0pt}}m{#1}}
\newcolumntype{C}[1]{>{\centering\let\newline\\\arraybackslash\hspace{0pt}}m{#1}}
\newcolumntype{R}[1]{>{\raggedleft\let\newline\\\arraybackslash\hspace{0pt}}m{#1}}

% special command for referencing RQs or table cells
\newcommand{\refRQ}[1]{\hfill {\tiny $\rightarrow RQ#1$}}
\newcommand{\refTab}[2]{{\scriptsize(#1,#2)}}

% rotate figures and tables
\usepackage{rotating}

% sub-figures
\usepackage{subcaption}

\usepackage[export]{adjustbox} % adds keys such as valign to \includegraphics (used for the logos on the title page)

% configure quote environment
\usepackage{quoting}
\quotingsetup{font={small}}
\quotingsetup{leftmargin=2em}
\quotingsetup{rightmargin=2em}
\quotingsetup{vskip=1.5ex}

% formatting for code 
\newcommand{\myCode}[1]{\texttt{\small#1}}

% pseudocode
\usepackage{algpseudocode}
\usepackage{algorithm}

% indentation of RQs
\newlength{\rqindent}\setlength{\rqindent}{\parindent+\widthof{RQ1\ }}

% appendix
\usepackage[titletoc]{appendix}

% colored boxes
\usepackage{tcolorbox}
\newtcolorbox{graybox}[1]{
    colback=black!3!white,
    colframe=black!15!white,
    coltitle=black!80!white,
    fonttitle=\sffamily\bfseries,
    toptitle=3pt,
    bottomtitle=3pt,
    title={#1}
}

% custom quotes
\usepackage{csquotes}
\newcommand*{\enq}[1]{\enquote{{\itshape#1}}}


%%%%%%%%%%%%%%%%%%%%%%%%%%%%%%%%%%%%%%%%

\newcommand{\thesistype}{Bachelor's Thesis} % Master's Thesis
\newcommand{\thesistitle}{Thesis Title}
\newcommand{\studentname}{Firstname Lastname}
\newcommand{\submissiondate}{\nth{1} Month Year}
\newcommand{\studentid}{12345678}
\newcommand{\examiner}{Prof. Dr. Sebastian Baltes}
\newcommand{\supervisor}{Title Firstname Lastname} % usually a doctoral or postdoctoral researcher

%%%%%%%%%%%%%%%%%%%%%%%%%%%%%%%%%%%%%%%%

\begin{document}
\hyphenchar\font=\string"7F % fix hyphenation of hyphenated words, see https://tex.stackexchange.com/a/63234/100364
\selectlanguage{english}

\frontmatter % roman numbers

\begin{titlepage}

	\begin{center}
		\includegraphics[height=1.3cm,valign=t]{figures/logos/logo-uhd.png}	
		\hfill
		\includegraphics[height=1.5cm,valign=t]{figures/logos/se-uhd-logo.png}	
	\end{center}
	
	\vspace{2cm}

    \begin{center}
        \large
        Heidelberg University\\
        Institute of Computer Science\\
        Software Engineering Group
    \end{center}

    \vspace{3cm}
    
	\begin{center}
		{\large \thesistype}\\
		\vspace{1.5\baselineskip}
		{\LARGE\textsf{\textbf{\thesistitle}}}\\
        \vspace{1.5\baselineskip}
        {\large \studentname}
	\end{center}

	\vfill

    \begin{center}
        \large
        \begin{tabular}{@{}ll@{}}
        \emph{Student Name:} & \studentname \\
        \emph{Student ID:} & \studentid \\
        \emph{Examiner:} & \examiner \\
        \emph{Topic Supervisor:} & \supervisor \\
        \emph{Submission Date:} & \submissiondate
        \end{tabular}
    \end{center}

    \vspace{2cm}

\end{titlepage}

\begin{otherlanguage}{ngerman}
\chapter*{Eigenständigkeitserklärung}
	
% https://www.hochschule-bayern.de/fileadmin/daten/Positionspapiere/KI-Leitlinie_Hochschule-Bayern-19052025.pdf
% https://www.informatik.uni-heidelberg.de/c/image/f/default/pdfs/BSc-template.pdf
Hiermit versichere ich, dass \dots

\begin{enumerate}
    \item ich die vorliegende Arbeit eigenständig verfasst und keine anderen als die angegebenen Quellen und Hilfsmittel verwendet habe.
    \item ich die Inhalte der Arbeit vollständig verstanden habe und selbstständig vertreten kann.
    \item ich die volle Verantwortung für die wissenschaftliche Qualität und den Inhalt der vorliegenden Arbeit, die gewählte Methodik und den Erstellungsprozess sowie die zitierte Literatur übernehme.
    \item die Arbeit bisher nicht anderweitig als Prüfungsleistung eingereicht wurde.
    \item ich alle wörtlich oder inhaltlich aus Quellen übernommenen oder mit Unterstützung von KI generierten Inhalte entsprechend \href{https://backend.uni-heidelberg.de/de/dokumente/satzung-zur-sicherung-guter-wissenschaftlicher-praxis-und-zum-umgang-mit-fehlverhalten-in-der/download}{anerkannten wissenschaftlichen Grundsätzen} bzw. den \href{https://backend.heiskills.uni-heidelberg.de/de/dokumente/ki-leitlinie-de/download}{geltenden Regelungen zur Kennzeichnung von KI-Inhalten} kenntlich gemacht habe.
    %Ausgenommen von der expliziten Kenntlichmachung sind orthografische oder grammatikalische Korrekturen, Übersetzungen sowie nicht-sinnverändernde Verbesserungen von Formulierungen.
    \item ich die Arbeit unter Einhaltung der \href{https://www.wissenschaftliche-integritaet.de/de/kodex}{guten wissenschaftlichen Praxis} erstellt habe.
    \item ich ausschließlich die unten genannten KI-Werkzeuge verwendet habe, deren Nutzung von dem Prüfer oder der Prüferin als Hilfsmittel zugelassen wurde.
    \item ich die unten genannten KI-Werkzeuge lediglich als Hilfsmittel genutzt habe, die von KI generierten Inhalte kritisch überprüft habe und mein eigenständiger kognitiver sowie kreativer Einfluss in dieser Arbeit überwiegt.
\end{enumerate}

Ich bin damit einverstanden, dass diese elektronische Fassung universitätsintern mithilfe von KI-Werkzeugen und Plagiatssoftware mit dem Zusatzmaterial der Arbeit abgeglichen sowie auf Plagiate, Inkonsistenzen, Fehler und halluzinierte Informationen überprüft wird.

\bigskip

Verwendete KI-Werkzeuge {\small(basierend auf dem \href{https://backend.heiskills.uni-heidelberg.de/de/dokumente/ki-formular-de/download}{Musterformular} der Universität Heidelberg)}:

% https://backend.heiskills.uni-heidelberg.de/de/dokumente/ki-formular-de/download
% https://backend.heiskills.uni-heidelberg.de/de/dokumente/ai-sample-form-en/download
\begin{center}
\small
\begin{tabularx}{\linewidth}{|l|X|X|}
\hline
\textbf{KI-Werkzeug} & \textbf{Verwendungszweck} & \textbf{Kommentare (optional)} \\
\hline
ChatGPT & Erstellung eines Gliederungsentwurfs. & -- \\
\hline
DeepL & Übersetzung von deutschen Texten ins Englische. & Es erfolgte eine Kontrolle durch einen Muttersprachler. \\ \hline
Claude Code & Unterstützung bei der Entwicklung des Tool-Prototyps. & Jeglicher generierte Code wurde von mir begutachtet und ggf. überarbeitet. \\
\hline
\end{tabularx} 
\end{center}

\vspace{2cm}
	
\noindent\begin{tabularx}{\textwidth}{XXXX}
	&&& \\
	\cline{1-1}\cline{4-4}\footnotesize\centering Ort und Datum		&&& \footnotesize\centering Unterschrift
\end{tabularx}
	
\newpage

% seminar-report: begin german-abstract (removed by convert-to-seminar-report.sh)
\chapter*{Zusammenfassung}

\todo{Write the German abstract.}

\remark{SB}{Die Zusammenfassung muss auf Deutsch und auf Englisch geschrieben werden. Die Zusammenfassung sollte zwischen einer halben und einer ganzen Seite lang sein. Sie soll den
Kontext der Arbeit, die Problemstellung, die Zielsetzung und die entwickelten Methoden
sowie Erkenntnisse bzw. Ergebnisse übersichtlich und verständlich beschreiben.}
% seminar-report: end german-abstract

\end{otherlanguage}

\chapter*{Abstract}

% seminar-report: begin abstract (replaced by convert-to-seminar-report.sh)
\todo{Write the abstract.}

\remark{SB}{The abstract has to be written in German and English. It should be between half a
page and one page in length. It should describe the context of the thesis, the problem
setting, the objectives, the methods developed in this thesis, and the key insights and
results in a clear and readable way.}
% seminar-report: end abstract

% seminar-report: begin acknowledgments (removed by convert-to-seminar-report.sh)
\chapter*{Acknowledgments}

\remark{SB}{This chapter is optional. You can thank the people who supported you during your thesis, such as your supervisors, study participants, or family and friends.}

\lipsum[1-4]
% seminar-report: end acknowledgments

{\hypersetup{linkcolor=black}\tableofcontents} % black entries in the table of contents

%\cleardoublepage
\clearpage

\mainmatter % arabic numbers

% https://se-uhd.de/q+a/general lists the \chapter headings below and the comments after them.
% You can delete these comments in your thesis.
% seminar-report: begin chapters (replaced by convert-to-seminar-report.sh)
\chapter{Introduction} % (including a problem statement)
\label{ch:introduction}

% you can use quotations at the start of each chapter using this environment
\begin{chapter-quotation}{Ward Cunningham, The Pragmatic Programmer (2000)}
If you don't think carefully, you might think that programming is just typing statements in a programming language.
\end{chapter-quotation}

\todo{Write the introduction.}

\remark{SB}{The introduction outlines the context, background, and motivation of your thesis. You can start from the introduction of your research proposal and update it based on what you learned during the thesis.}

\remark{SB}{Start with the broad context, narrow it down to the problem that your thesis addresses, and argue why the problem matters. Refer to the most important related work and explain the knowledge gap that your thesis closes. Structure the thesis around its outcome rather than the chronological order of your work. After the problem statement and the research questions, summarize your contributions and main results and give an overview of the remaining chapters.}

% example paragraph with citations
The software engineering research community is often criticized for not addressing the problems that practitioners actually face during their work~\cite{Briand2012}.
At the same time, software developers base their beliefs on personal experience rather than on empirical findings~\cite{DevanbuZimmermannOthers2016}.
We summarize this discussion in \autoref{ch:related-work}.

\begin{java}
public class HelloWorld {
  public static void main(String[] args) {
    System.out.println("Hello, World!");
  }
}
\end{java}

To embed quotations, such as the statement attributed to W. Edwards Deming that \enq{without data, you're just another person with an opinion}, use the custom \texttt{\textbackslash enq} command.

\todo{Formulate the problem statement.}

\remark{SB}{State the problem that your thesis addresses in a few sentences. The research questions should follow from it.}

\begin{graybox}{Problem Statement}
Problem statement.
\end{graybox}

\todo{Formulate the research question(s).}

\remark{SB}{Formulate the objectives of your thesis as research questions. They usually refine the questions from your research proposal. Answer each research question in \autoref{ch:results}.}

\begin{graybox}{Research Questions}
\begin{description}[style=multiline, labelindent=\parindent, leftmargin=\rqindent, itemsep=-1ex]
\item[RQ1] ...?
\item[RQ2] ...?
\item[RQ3] ...?
\end{description}
\end{graybox}

\chapter{Background} % (can sometimes be integrated into the introduction)
\label{ch:background}

\todo{Write the background chapter.}

\remark{SB}{Introduce the concepts, terminology, and technologies that readers need to follow the rest of the thesis. Cover only what later chapters use, and discuss prior research on the same or a similar problem in \autoref{ch:related-work}.}

\lipsum[1-12]

\chapter{Related Work}
\label{ch:related-work}

\todo{Write the related work chapter.}

\remark{SB}{Discuss prior research on the problem that your thesis addresses and explain how your work differs from it or builds on it. Group the work by topic or approach instead of summarizing one paper after another. Most papers only need a sentence or two. The related work section of your research proposal is a starting point.}

\lipsum[1-4]

\chapter{Research Method} % (including a data availability statement)
\label{ch:research-method}

\todo{Describe your research method.}

\remark{SB}{Describe how you answer each research question, including the research methods, data collection, participant recruitment, analysis procedures, and statistical methods. State which research question each step answers, and give enough detail for others to replicate your study.}

\remark{SB}{If you built a tool prototype or a pipeline, describe it in \autoref{ch:design} and use this chapter to explain how you evaluate it. For a systematic literature review, name the review guidelines that you follow here and describe the concrete review pipeline in \autoref{ch:design}.}

\lipsum[1-4]

\section{Data Availability}
\label{sec:data-availability}

\todo{Replace NAME in the URL with the name of your repository.}

\remark{SB}{If the dataset is too large to host on GitHub, you can publish the online appendix on Zenodo and reference the Zenodo dataset (example: \url{https://zenodo.org/records/15130575}). In that case, still link the GitHub repository with the source code and version history.}

We provide all source code, data, and study materials as part of an online appendix hosted on GitHub.\footnote{\url{https://github.com/se-uhd/thesis-NAME}}

\chapter{Design (and Implementation)}
\label{ch:design}

\todo{Describe the design and, if applicable, the implementation of your tool or pipeline.}

\remark{SB}{Describe what you built or set up to answer your research questions, such as a tool prototype, a data analysis pipeline, software for an online experiment, or the review pipeline of a systematic literature review, from the search to data extraction.}

\remark{SB}{Most students implement a tool, a data analysis pipeline, or scripts that automate the search or screening of a systematic literature review. If you did, remove only the parentheses from the chapter title. If you did not, for example because you searched and screened the literature manually, title the chapter ``Design''.}

\remark{SB}{Start with an overview of the design, for example an architecture diagram or a figure of the pipeline stages, and explain the main design decisions and the alternatives that you considered. If you implemented the design, describe how you implemented each component or stage, including the technologies, scripts, or tools that you used. Implementation details that do not help readers understand or reuse your work belong in the online appendix (see \autoref{sec:data-availability}). Remove this chapter if your thesis includes no tool or pipeline.}

\lipsum[1-4]

\chapter{Results}
\label{ch:results}

\todo{Report your results.}

\remark{SB}{Report the results for each research question in the order of the questions. Use tables and figures where they make results easier to compare, and leave the interpretation to \autoref{ch:discussion}.}

\lipsum[1-4]

\chapter{Discussion} % (including a critical reflection)
\label{ch:discussion}

\todo{Discuss your results.}

\remark{SB}{Interpret your results, relate them to the related work, and explain what they mean for researchers and practitioners. Discuss the threats to the validity of your study and how you addressed them. Do not introduce new results in this chapter.}

\lipsum[1-4]

\chapter{Conclusion and Future Work} % (including an outlook)
\label{ch:conclusion}

\todo{Write the conclusion.}

\remark{SB}{Summarize the answers to your research questions and your contributions, and outline directions for future work. Do not introduce new results in this chapter.}

\lipsum[1-4]

% seminar-report: end chapters

{\hypersetup{linkcolor=black}\listoffigures} % black entries in the list of figures

{\hypersetup{linkcolor=black}\listoftables} % black entries in the list of tables

\bibliographystyle{plainnat}
\bibliography{literature} % References

\begin{appendices}

\chapter{Title}

\remark{SB}{Use the appendix for material that would interrupt the main chapters, such as interview guides, questionnaires, or additional tables. Remove it if you do not need it.}

% seminar-report: begin online-appendix (removed by convert-to-seminar-report.sh)
\remark{SB}{Provide all source code, data, and study materials in the online appendix (see \autoref{sec:data-availability}), even if you include parts of them here.}
% seminar-report: end online-appendix

\end{appendices}

\backmatter

\end{document}
